% ============================================================
%  GUÍA 3: FUNCIONES LINEALES
%  Curso de Introducción 2do Año
%  Autor: Ing. Gabriel Astudillo
%  Clases cubiertas: 9 a 12
% ============================================================

\documentclass[12pt, a4paper]{article}

% ── Paquetes ─────────────────────────────────────────────────

\usepackage{mathwizards-guia}
\mwguia{Guía 3 — Funciones Lineales}{Ing. Gabriel Astudillo $\cdot$ 2do Año}

\begin{document}
% ============================================================

% ── Portada / Título ─────────────────────────────────────────
\mwportada{Guía de Estudio 3}{Funciones Lineales}{Función Lineal y Afín, Plano Cartesiano, Gráficas y Pendiente}

\vspace{4pt}
\begin{cuadroinfo}
  \small
  \textbf{Presentación:} Las funciones lineales describen muchísimas situaciones cotidianas: el costo de una compra, la distancia recorrida a velocidad constante, el saldo de una cuenta. En esta guía aprenderás a distinguir una función lineal de una afín, a ubicar puntos en el plano cartesiano, a graficar rectas por tabulación y por puntos de corte, a analizar la pendiente y a modelar problemas con funciones de costo. Cada sección incluye \textbf{ejercicios resueltos} y \textbf{ejercicios propuestos}.
  \\[4pt]
  \textbf{Nivel de Dificultad:}\quad
  \facil{} Básico / Directo \quad
  \medio{} Intermedio \quad
  \dificil{} Desafiante
\end{cuadroinfo}

\vspace{8pt}

% ============================================================
\section{Conceptos Básicos: Función Lineal y Función Afín}
% ============================================================

\subsection{Definiciones}

\begin{cuadroteoria}
  \textbf{Función:} una regla que a cada valor de la \textbf{variable independiente} $x$ le asigna exactamente un valor de la \textbf{variable dependiente} $y$ (o $F(x)$). La variable $y$ ``depende'' de $x$.

  \textbf{Función lineal (proporcionalidad directa):}
  $$F(x) = mx \qquad \text{(su gráfica pasa por el origen)}$$
  Al duplicar $x$, se duplica $F(x)$: las variables son \emph{directamente proporcionales}.

  \textbf{Función afín:}
  $$F(x) = mx + b \qquad \text{(con } b \neq 0 \text{; no pasa por el origen)}$$

  \textbf{Función constante:}
  $$F(x) = b \qquad \text{(pendiente cero; recta horizontal)}$$
\end{cuadroteoria}

\begin{cuadroteoria}
  \textbf{Forma general:} $F(x) = mx + b$, donde:
  \begin{itemize}[leftmargin=*]
    \item $m$ es la \textbf{pendiente}: indica cuánto cambia $y$ por cada unidad que cambia $x$.
    \item $b$ es el \textbf{coeficiente de posición} (intercepto con el eje $y$): el valor de $F(0)$.
  \end{itemize}
\end{cuadroteoria}

\begin{cuadroadvertencia}
  \textbf{No confundir:}
  \begin{itemize}[leftmargin=*]
    \item \emph{Lineal} ($F(x) = mx$): pasa por el origen, proporcionalidad directa.
    \item \emph{Afín} ($F(x) = mx + b$ con $b \neq 0$): no pasa por el origen.
    \item En muchos textos escolares a ambas se les llama ``funciones lineales''; aquí usaremos la distinción de arriba.
  \end{itemize}
\end{cuadroadvertencia}

\subsection{Ejercicios resueltos}

\textbf{Ejemplo R1.} Clasifica cada función en lineal (proporcionalidad directa), afín o constante:
$$a) \; F(x) = 5x \qquad b) \; F(x) = 3x - 2 \qquad c) \; F(x) = -4 \qquad d) \; F(x) = \frac{2}{3}x$$

\begin{cuadrosolucion}
  \textbf{Solución:}
  \begin{itemize}
    \item $a) \; F(x) = 5x$: \textbf{lineal} ($m = 5$, $b = 0$, pasa por el origen).
    \item $b) \; F(x) = 3x - 2$: \textbf{afín} ($m = 3$, $b = -2 \neq 0$).
    \item $c) \; F(x) = -4$: \textbf{constante} (pendiente $m = 0$).
    \item $d) \; F(x) = \dfrac{2}{3}x$: \textbf{lineal} ($m = \frac{2}{3}$, $b = 0$).
  \end{itemize}
\end{cuadrosolucion}

\textbf{Ejemplo R2.} En la función $F(x) = -2x + 7$, identifica la pendiente, el coeficiente de posición y calcula $F(3)$.

\begin{cuadrosolucion}
  \textbf{Solución:} Comparando con $F(x) = mx + b$:
  \begin{itemize}
    \item Pendiente: $m = -2$.
    \item Coeficiente de posición: $b = 7$.
    \item $F(3) = -2 \cdot 3 + 7 = -6 + 7 = 1$.
  \end{itemize}
\end{cuadrosolucion}

\textbf{Ejemplo R3.} Un auto recorre 80 km por hora. Escribe la función que da la distancia recorrida en función del tiempo y calcula cuánto avanza en 2,5 horas.

\begin{cuadrosolucion}
  \textbf{Solución:} La distancia es proporcional al tiempo: $F(t) = 80t$ (lineal, pasa por el origen).
  $$F(2{,}5) = 80 \cdot 2{,}5 = 200\ \text{km}.$$
\end{cuadrosolucion}

\newpage

\subsection{Ejercicios propuestos}

\begin{enumerate}[leftmargin=*]
  \item \facil{} Clasifica en lineal, afín o constante: $a) \; F(x) = 4x + 1 \qquad b) \; F(x) = \dfrac{x}{2} \qquad c) \; F(x) = 9 \qquad d) \; F(x) = -3x$

  \item \facil{} Identifica $m$ y $b$ en cada función: $a) \; F(x) = 6x - 1 \qquad b) \; F(x) = -x + 4 \qquad c) \; F(x) = \dfrac{3}{4}x$

  \item \facil{} Calcula: si $F(x) = 2x + 5$, halla $F(0)$, $F(1)$ y $F(-2)$.

  \item \facil{} Calcula: si $G(x) = -3x + 1$, halla $G(2)$ y $G(-1)$.

  \item \medio{} ¿Cuál es la pendiente de una función constante? ¿Por dónde pasa la gráfica de $F(x) = mx$?

  \item \medio{} Un celular cuesta \$12\,000 y su valor baja \$800 por cada año que pasa. Escribe la función $V(t)$ del valor en función de los años $t$, e indica si es lineal o afín.

  \item \dificil{} Si $F(x) = mx + b$, $F(1) = 5$ y $F(3) = 11$, determina $m$ y $b$. (Pista: plantea un sistema de dos ecuaciones.)
\end{enumerate}

%\newpage %ajuste pag

% ============================================================
\section{Plano Cartesiano y Coordenadas}
% ============================================================

\subsection{Estructura del plano cartesiano}

\begin{cuadroteoria}
  El \textbf{plano cartesiano} está formado por dos rectas numéricas perpendiculares:
  \begin{itemize}[leftmargin=*]
    \item \textbf{Eje $X$} (eje de las abscisas): recta horizontal.
    \item \textbf{Eje $Y$} (eje de las ordenadas): recta vertical.
    \item Se cortan en el \textbf{origen} $(0, 0)$.
    \item Dividen el plano en cuatro \textbf{cuadrantes}: I (arriba a la derecha), II (arriba a la izquierda), III (abajo a la izquierda), IV (abajo a la derecha).
  \end{itemize}

  \textbf{Par ordenado} $(x, y)$: la primera coordenada $x$ indica el desplazamiento horizontal y la segunda $y$ el vertical.
\end{cuadroteoria}

\begin{cuadroteoria}
  \textbf{Puntos sobre los ejes:}
  \begin{itemize}[leftmargin=*]
    \item Sobre el eje $X$: $y = 0$, se escriben $(x, 0)$.
    \item Sobre el eje $Y$: $x = 0$, se escriben $(0, y)$.
  \end{itemize}

  \textbf{Coordenadas fraccionarias:} se ubican de forma aproximada entre dos enteros. Por ejemplo, $\left(-\dfrac{9}{4}, -\dfrac{3}{2}\right) = (-2{,}25, -1{,}5)$: un poco a la izquierda de $-2$ y un poco abajo de $-1$.
\end{cuadroteoria}

\subsection{Ejercicios resueltos}

\textbf{Ejemplo R1.} Ubica los puntos $A(3, 2)$, $B(-2, 4)$, $C(-1, -3)$, $D(2, -1)$ y $E(0, 3)$ en el plano cartesiano e indica en qué cuadrante está cada uno.

\begin{cuadrosolucion}
  \textbf{Solución:}
  \begin{center}
    \begin{tikzpicture}[scale=0.75]
      \draw[thin, gray!40] (-3.5,-3.5) grid (4.5,4.5);
      \draw[thick, -Latex] (-3.8,0) -- (4.8,0) node[right] {$x$};
      \draw[thick, -Latex] (0,-3.8) -- (0,4.8) node[above] {$y$};
      \draw[gray] (0,0) node[below left] {$0$};
      \filldraw[mwprimario] (3,2) circle (2.5pt) node[above right] {$A(3,2)$};
      \filldraw[mwprimario] (-2,4) circle (2.5pt) node[above] {$B(-2,4)$};
      \filldraw[mwprimario] (-1,-3) circle (2.5pt) node[below] {$C(-1,-3)$};
      \filldraw[mwprimario] (2,-1) circle (2.5pt) node[below right] {$D(2,-1)$};
      \filldraw[mwprimario] (0,3) circle (2.5pt) node[right] {$E(0,3)$};
    \end{tikzpicture}
  \end{center}
  \begin{itemize}
    \item $A(3, 2)$: cuadrante I.
    \item $B(-2, 4)$: cuadrante II.
    \item $C(-1, -3)$: cuadrante III.
    \item $D(2, -1)$: cuadrante IV.
    \item $E(0, 3)$: sobre el eje $Y$ (no está en ningún cuadrante).
  \end{itemize}
\end{cuadrosolucion}

\newpage

\textbf{Ejemplo R2.} Ubica los puntos con coordenadas fraccionarias $P\left(-\dfrac{9}{4}, -\dfrac{3}{2}\right)$ y $Q\left(\dfrac{5}{2}, \dfrac{1}{2}\right)$.

\begin{cuadrosolucion}
  \textbf{Solución:} Convertimos a decimales:
  $$P\left(-\frac{9}{4}, -\frac{3}{2}\right) = P(-2{,}25, -1{,}5) \qquad Q\left(\frac{5}{2}, \frac{1}{2}\right) = Q(2{,}5, 0{,}5).$$
  \begin{center}
    \begin{tikzpicture}[scale=0.75]
      \draw[thin, gray!40] (-3.5,-2.5) grid (3.5,2.5);
      \draw[thick, -Latex] (-3.8,0) -- (3.8,0) node[right] {$x$};
      \draw[thick, -Latex] (0,-2.8) -- (0,2.8) node[above] {$y$};
      \draw[gray] (0,0) node[below left] {$0$};
      \filldraw[mwprimario] (-2.25,-1.5) circle (2.5pt) node[below left] {$P$};
      \filldraw[mwprimario] (2.5,0.5) circle (2.5pt) node[above right] {$Q$};
    \end{tikzpicture}
  \end{center}
  $P$ queda un poco a la izquierda de $-2$ y entre $-1$ y $-2$ hacia abajo; $Q$ queda entre $2$ y $3$ en $x$, y a media unidad arriba del eje $X$.
\end{cuadrosolucion}

\subsection{Ejercicios propuestos}

\begin{enumerate}[leftmargin=*, resume]
  \item \facil{} Indica el cuadrante (o el eje) de cada punto: $a)\ (5, 3)$ \quad $b)\ (-1, 6)$ \quad $c)\ (-4, -2)$ \quad $d)\ (0, -5)$ \quad $e)\ (7, 0)$

  \item \facil{} Escribe las coordenadas del punto: $a) \; 3$ unidades a la derecha y $2$ arriba del origen. $b) \; 4$ unidades a la izquierda y $1$ abajo.

  \item \facil{} ¿En qué cuadrante está un punto con $x < 0$ e $y < 0$? ¿Y con $x > 0$ e $y < 0$?

  \item \medio{} Ubica en el plano cartesiano: $a)\ (-1, -3)$ \quad $b)\ (2, -8)$ \quad $c)\ \left(-\dfrac{9}{4}, -\dfrac{3}{2}\right)$ \quad $d)\ (0, 2)$ \\
        $e)\ (-5, 3)$ \quad $f)\ (-5, 0)$ \quad $g)\ (0, -3)$ \quad $h)\ \left(-6, -\dfrac{3}{2}\right)$

  \item \medio{} Escribe el cuadrante en el que se ubica cada punto del ejercicio anterior.

  \item \dificil{} Un punto está en el cuadrante II y sus coordenadas son números enteros cuya suma es $-1$. ¿Cuáles pueden ser sus coordenadas? Encuentra dos posibilidades.
\end{enumerate}

\newpage

% ============================================================
\section{Identificación del Tipo de Función}
% ============================================================

\subsection{A partir de la expresión algebraica}

\begin{cuadroteoria}
  Dada $F(x) = mx + b$:
  \begin{itemize}[leftmargin=*]
    \item $b = 0$ $\Rightarrow$ \textbf{lineal} (proporcionalidad directa): pasa por el origen.
    \item $b \neq 0$ $\Rightarrow$ \textbf{afín}: no pasa por el origen.
    \item $m = 0$ $\Rightarrow$ \textbf{constante}: recta horizontal.
  \end{itemize}
\end{cuadroteoria}

\subsection{A partir de dos puntos}

\begin{cuadroteoria}
  Dados dos puntos de la función:
  \begin{itemize}[leftmargin=*]
    \item Si al aumentar $x$ aumenta $y$ $\Rightarrow$ \textbf{creciente} (pendiente positiva).
    \item Si al aumentar $x$ disminuye $y$ $\Rightarrow$ \textbf{decreciente} (pendiente negativa).
    \item Si ambos puntos tienen la misma ordenada $y$ $\Rightarrow$ \textbf{constante} (pendiente cero).
  \end{itemize}
\end{cuadroteoria}

\subsection{Ejercicios resueltos}

\textbf{Ejemplo R1.} Clasifica según su expresión: $a) \; F(x) = 7x \qquad b) \; F(x) = -x + 3 \qquad c) \; F(x) = \dfrac{2}{5}x \qquad d) \; F(x) = 12$

\begin{cuadrosolucion}
  \textbf{Solución:}
  \begin{itemize}
    \item $a) \; F(x) = 7x$: lineal (pasa por el origen).
    \item $b) \; F(x) = -x + 3$: afín ($b = 3 \neq 0$).
    \item $c) \; F(x) = \dfrac{2}{5}x$: lineal (proporcionalidad directa).
    \item $d) \; F(x) = 12$: constante ($m = 0$).
  \end{itemize}
\end{cuadrosolucion}

\textbf{Ejemplo R2.} La función pasa por los puntos $(1, 3)$ y $(-1, 3)$. ¿Es creciente, decreciente o constante?

\begin{cuadrosolucion}
  \textbf{Solución:} Ambos puntos tienen la misma ordenada ($y = 3$). Por lo tanto, la función es \textbf{constante}: su gráfica es una recta horizontal $y = 3$.
\end{cuadrosolucion}

\textbf{Ejemplo R3.} Determina si la función que pasa por $(2, 5)$ y $(4, 9)$ es creciente o decreciente.

\begin{cuadrosolucion}
  \textbf{Solución:} Al pasar de $x = 2$ a $x = 4$, la imagen pasa de $5$ a $9$: aumentó. La función es \textbf{creciente}.
\end{cuadrosolucion}

\newpage

\subsection{Ejercicios propuestos}

\begin{enumerate}[leftmargin=*, resume]
  \item \facil{} Clasifica en lineal, afín o constante: $a) \; F(x) = 3x - 8 \qquad b) \; F(x) = -5x \qquad c) \; F(x) = 0x + 6 \qquad d) \; F(x) = \dfrac{x}{7}$

  \item \facil{} Determina si la función que pasa por los puntos dados es creciente, decreciente o constante: $a) \; (0, 2)$ y $(3, 2)$ \quad $b) \; (1, 1)$ y $(2, 4)$ \quad $c) \; (5, 8)$ y $(6, 3)$

  \item \medio{} La función pasa por $(-2, 4)$ y $(1, 4)$. ¿Qué puedes afirmar sobre su pendiente?

  \item \medio{} Escribe la forma general de una función afín que sea decreciente y cuyo coeficiente de posición sea $5$. (Hay infinitas; da un ejemplo.)

  \item \dificil{} Si $F(x) = mx + b$ pasa por los puntos $(1, 2)$ y $(3, 8)$, calcula $m$ y $b$. (Pista: la pendiente es $\dfrac{\text{cambio en } y}{\text{cambio en } x}$.)
\end{enumerate}

%\newpage %ajuste pag

% ============================================================
\section{Gráfica de Funciones Lineales}
% ============================================================

\subsection{Método 1: Tabulación}

\begin{cuadropasos}
  \textbf{Pasos:}
  \begin{enumerate}[leftmargin=*]
    \item Asigna valores a $x$ (por ejemplo, $-2, -1, 0, 1, 2$).
    \item Calcula $F(x)$ para cada valor.
    \item Completa la \textbf{tabla de valores}.
    \item Ubica los pares ordenados en el plano y únelos con una recta.
  \end{enumerate}
\end{cuadropasos}

\subsection{Método 2: Puntos de corte}

\begin{cuadropasos}
  \textbf{Pasos:}
  \begin{enumerate}[leftmargin=*]
    \item \textbf{Corte con el eje $Y$:} se hace $x = 0$; el punto es $(0, F(0)) = (0, b)$.
    \item \textbf{Corte con el eje $X$:} se resuelve $F(x) = 0$; el punto es $(x, 0)$.
    \item Se traza la recta que une ambos puntos.
  \end{enumerate}
\end{cuadropasos}

\begin{cuadroadvertencia}
  \textbf{Caso especial:} si la función es \emph{constante} ($F(x) = b$), la recta es horizontal y no corta el eje $X$ (a menos que $b = 0$).
\end{cuadroadvertencia}

\subsection{Ejercicios resueltos}

\textbf{Ejemplo R1.} Grafica $F(x) = 2x + 1$ por tabulación e indica el tipo de función.

\begin{cuadrosolucion}
  \textbf{Solución:}
  \begin{center}
    \begin{tabular}{c|ccccc}
      \toprule
      $x$             & $-2$ & $-1$ & $0$ & $1$ & $2$ \\
      \midrule
      $F(x) = 2x + 1$ & $-3$ & $-1$ & $1$ & $3$ & $5$ \\
      \bottomrule
    \end{tabular}
  \end{center}
  \begin{center}
    \begin{tikzpicture}[scale=0.75]
      \draw[thin, gray!40] (-3.5,-4) grid (3.5,6);
      \draw[thick, -Latex] (-3.8,0) -- (3.8,0) node[right] {$x$};
      \draw[thick, -Latex] (0,-4.2) -- (0,6.2) node[above] {$y$};
      \draw[gray] (0,0) node[below left] {$0$};
      \draw[domain=-2.2:2.2, smooth, very thick, mwprimario] plot (\x, {2*\x + 1});
      \filldraw[mwprimario] (0,1) circle (2.5pt) node[above right] {$(0,1)$};
      \filldraw[mwprimario] (2,5) circle (2.5pt) node[above right] {$(2,5)$};
    \end{tikzpicture}
  \end{center}
  Como $F(x) = 2x + 1$ tiene $b = 1 \neq 0$, es una función \textbf{afín} creciente.
\end{cuadrosolucion}

\newpage

\textbf{Ejemplo R2.} Grafica $F(x) = 3x - 6$ por el método de los puntos de corte e identifica el tipo de función.

\begin{cuadrosolucion}
  \textbf{Solución:}
  \begin{itemize}
    \item \textbf{Corte con el eje $Y$:} $F(0) = 3 \cdot 0 - 6 = -6$ $\Rightarrow$ punto $(0, -6)$.
    \item \textbf{Corte con el eje $X$:} $3x - 6 = 0 \Rightarrow 3x = 6 \Rightarrow x = 2$ $\Rightarrow$ punto $(2, 0)$.
  \end{itemize}
  \begin{center}
    \begin{tikzpicture}[scale=0.75]
      \draw[thin, gray!40] (-3.5,-7) grid (3.5,4);
      \draw[thick, -Latex] (-3.8,0) -- (3.8,0) node[right] {$x$};
      \draw[thick, -Latex] (0,-7.2) -- (0,4.2) node[above] {$y$};
      \draw[gray] (0,0) node[below left] {$0$};
      \draw[domain=-1.5:2.7, smooth, very thick, mwprimario] plot (\x, {3*\x - 6});
      \filldraw[mwprimario] (0,-6) circle (2.5pt) node[below left] {$(0,-6)$};
      \filldraw[mwprimario] (2,0) circle (2.5pt) node[below right] {$(2,0)$};
    \end{tikzpicture}
  \end{center}
  Es una función \textbf{afín} creciente ($m = 3 > 0$, $b = -6 \neq 0$).
\end{cuadrosolucion}

\newpage %ajuste pag

\textbf{Ejemplo R3.} Grafica $F(x) = -\dfrac{4}{3}x + 4$ por el método de los puntos de corte.

\begin{cuadrosolucion}
  \textbf{Solución:}
  \begin{itemize}
    \item \textbf{Corte con el eje $Y$:} $F(0) = 4$ $\Rightarrow$ punto $(0, 4)$.
    \item \textbf{Corte con el eje $X$:} $-\dfrac{4}{3}x + 4 = 0 \Rightarrow -\dfrac{4}{3}x = -4 \Rightarrow x = 3$ $\Rightarrow$ punto $(3, 0)$.
  \end{itemize}
  \begin{center}
    \begin{tikzpicture}[scale=0.75]
      \draw[thin, gray!40] (-2.5,-2) grid (4.5,5);
      \draw[thick, -Latex] (-2.8,0) -- (4.8,0) node[right] {$x$};
      \draw[thick, -Latex] (0,-2.2) -- (0,5.2) node[above] {$y$};
      \draw[gray] (0,0) node[below left] {$0$};
      \draw[domain=-0.5:3.6, smooth, very thick, mwprimario] plot (\x, {-4/3*\x + 4});
      \filldraw[mwprimario] (0,4) circle (2.5pt) node[above left] {$(0,4)$};
      \filldraw[mwprimario] (3,0) circle (2.5pt) node[below right] {$(3,0)$};
    \end{tikzpicture}
  \end{center}
  Es una función \textbf{afín} decreciente ($m = -\frac{4}{3} < 0$).
\end{cuadrosolucion}

\newpage

\textbf{Ejemplo R4.} Grafica $F(x) = -3x$ e identifica el tipo de función.

\begin{cuadrosolucion}
  \textbf{Solución:} Como $b = 0$, es una función \textbf{lineal} (proporcionalidad directa) y su gráfica pasa por el origen. Los dos cortes coinciden en $(0, 0)$, así que necesitamos un punto extra: con $x = 1$, $F(1) = -3$, punto $(1, -3)$.
  \begin{center}
    \begin{tikzpicture}[scale=0.75]
      \draw[thin, gray!40] (-3.5,-4) grid (3.5,3);
      \draw[thick, -Latex] (-3.8,0) -- (3.8,0) node[right] {$x$};
      \draw[thick, -Latex] (0,-4.2) -- (0,3.2) node[above] {$y$};
      \draw[gray] (0,0) node[below left] {$0$};
      \draw[domain=-1.3:1.3, smooth, very thick, mwprimario] plot (\x, {-3*\x});
      \filldraw[mwprimario] (0,0) circle (2.5pt) node[below left] {$(0,0)$};
      \filldraw[mwprimario] (1,-3) circle (2.5pt) node[below right] {$(1,-3)$};
    \end{tikzpicture}
  \end{center}
\end{cuadrosolucion}

\subsection{Ejercicios propuestos}

\begin{enumerate}[leftmargin=*, resume]
  \item \facil{} Grafica por tabulación: $F(x) = x + 3$ (usa $x = -2, -1, 0, 1, 2$).

  \item \facil{} Grafica por tabulación: $F(x) = -2x + 1$.

  \item \medio{} Grafica por el método de los puntos de corte: $a) \; F(x) = 3x - 6 \qquad b) \; F(x) = -2x + 8$

  \item \medio{} Grafica por el método de los puntos de corte: $a) \; F(x) = -\dfrac{4}{3}x + 4 \qquad b) \; F(x) = \dfrac{2}{3}x - 6$

  \item \medio{} Grafica: $a) \; F(x) = -3x \qquad b) \; F(x) = 0x - 6$ e identifica el tipo de cada una.

  \item \dificil{} Sin graficar, determina si la función corta el eje $X$ en un valor positivo o negativo: $F(x) = \dfrac{2}{3}x - 6$. Justifica.
\end{enumerate}

\newpage

% ============================================================
\section{Análisis de la Pendiente}
% ============================================================

\subsection{Significado de la pendiente}

\begin{cuadroteoria}
  La \textbf{pendiente} $m$ de $F(x) = mx + b$ indica la inclinación de la recta:
  \begin{itemize}[leftmargin=*]
    \item $m > 0$: función \textbf{creciente}; la recta asciende de izquierda a derecha.
    \item $m < 0$: función \textbf{decreciente}; la recta desciende de izquierda a derecha.
    \item $m = 0$: función \textbf{constante}; recta horizontal paralela al eje $X$.
  \end{itemize}

  \textbf{Cálculo de la pendiente entre dos puntos} $(x_1, y_1)$ y $(x_2, y_2)$:
  $$m = \frac{y_2 - y_1}{x_2 - x_1}$$
\end{cuadroteoria}

\begin{cuadroextra}
  \textbf{Interpretación práctica:} en una función de costo, la pendiente es el \emph{costo por unidad} (precio unitario). En una función de distancia-tiempo, la pendiente es la \emph{velocidad}.
\end{cuadroextra}

\subsection{Ejercicios resueltos}

\textbf{Ejemplo R1.} Calcula la pendiente de la recta que pasa por $(1, 2)$ y $(4, 11)$ e indica si la función es creciente o decreciente.

\begin{cuadrosolucion}
  \textbf{Solución:}
  $$m = \frac{11 - 2}{4 - 1} = \frac{9}{3} = 3.$$
  Como $m = 3 > 0$, la función es \textbf{creciente}.
\end{cuadrosolucion}

\textbf{Ejemplo R2.} Calcula la pendiente de la recta que pasa por $(3, 8)$ y $(5, 2)$.

\begin{cuadrosolucion}
  \textbf{Solución:}
  $$m = \frac{2 - 8}{5 - 3} = \frac{-6}{2} = -3.$$
  Como $m = -3 < 0$, la función es \textbf{decreciente}.
\end{cuadrosolucion}

\textbf{Ejemplo R3.} La gráfica de una función es una recta horizontal que pasa por el punto $(2, 5)$. ¿Cuál es su pendiente y cuál es su ecuación?

\begin{cuadrosolucion}
  \textbf{Solución:} Recta horizontal $\Rightarrow$ pendiente $m = 0$ y todos los puntos tienen $y = 5$. Su ecuación es $F(x) = 5$ (función constante).
\end{cuadrosolucion}

\newpage

\subsection{Ejercicios propuestos}

\begin{enumerate}[leftmargin=*, resume]
  \item \facil{} Indica si cada función es creciente, decreciente o constante según su pendiente: $a)\ F(x) = 5x - 1$ \quad $b)\ F(x) = -\dfrac{1}{2}x + 3$ \quad $c)\ F(x) = -4x$ \quad $d)\ F(x) = 2$

  \item \facil{} Calcula la pendiente entre los pares de puntos: $a)\ (0, 0)$ y $(2, 6)$ \enskip $b)\ (1, 4)$ y $(3, 10)$ \enskip $c)\ (2, 7)$ y $(5, 1)$

  \item \medio{} Calcula la pendiente entre: $a) \; (-1, 3)$ y $(2, -3)$ \quad $b) \; (4, 1)$ y $(4, 9)$ (¿qué observas?)

  \item \medio{} Determina la pendiente de la recta $F(x) = \dfrac{2}{3}x - 6$ y di si su gráfica asciende o desciende.

  \item \dificil{} Escribe una función afín con pendiente $m = -2$ que pase por el punto $(1, 4)$.

  \item \dificil{} Una recta pasa por los puntos $(a, 5)$ y $(a + 2, 9)$. Calcula su pendiente sin conocer el valor de $a$.
\end{enumerate}

% \newpage %ajuste pag

% ============================================================
\section{Problemas de Aplicación (Modelado)}
% ============================================================

\subsection{El método del modelado}

\begin{cuadropasos}
  \textbf{Pasos para modelar un problema con una función lineal:}
  \begin{enumerate}[leftmargin=*]
    \item \textbf{Extrae los datos} del enunciado: identifica las variables y los valores dados.
    \item \textbf{Calcula la constante de proporcionalidad} (precio unitario, razón de cambio): divide el costo total entre la cantidad.
    \item \textbf{Plantea la función:} $F(x) = mx$ si no hay costo fijo, o $F(x) = mx + b$ si hay un costo fijo inicial.
    \item \textbf{Verifica} con los valores dados en el enunciado.
  \end{enumerate}
\end{cuadropasos}

\subsection{Ejercicios resueltos}

\textbf{Ejemplo R1.} He comprado 1 kilo y medio de tomates y me han costado \$1\,200. Escribe la función que da el costo de los tomates en función de su peso.

\begin{cuadrosolucion}
  \textbf{Solución:}
  \begin{enumerate}
    \item Datos: $1{,}5$ kg cuestan \$1\,200.
    \item Constante de proporcionalidad (precio por kilo): $m = \dfrac{1200}{1{,}5} = 800$ \$/kg.
    \item Como no hay costo fijo, la función es lineal: $F(x) = 800x$.
    \item Verificación: $F(1{,}5) = 800 \cdot 1{,}5 = 1\,200$ $\checkmark$
  \end{enumerate}
\end{cuadrosolucion}

\textbf{Ejemplo R2.} Una empresa de transporte cobra \$1\,500 por el primer kilómetro y \$500 por cada kilómetro adicional. Escribe la función del costo en función de los kilómetros recorridos y calcula el costo de un viaje de 8 km.

\begin{cuadrosolucion}
  \textbf{Solución:}
  \begin{enumerate}
    \item Hay un costo fijo inicial (por el primer kilómetro) y un costo por kilómetro adicional.
    \item Planteamos: si $x$ es el número de kilómetros, el costo es $F(x) = 1\,500 + 500(x - 1)$ (el primer km ya está incluido).
    \item Simplificamos: $F(x) = 1\,500 + 500x - 500 = 500x + 1\,000$.
    \item Para $x = 8$: $F(8) = 500 \cdot 8 + 1\,000 = 4\,000 + 1\,000 = 5\,000$.
  \end{enumerate}
  El viaje de 8 km cuesta \$5\,000. Nota: aquí la función es \emph{afín} porque existe un costo fijo.
\end{cuadrosolucion}

\textbf{Ejemplo R3.} Un estacionamiento cobra \$2\,000 fijos más \$300 por cada hora. Escribe la función del costo y determina cuántas horas se pueden estacionar con \$5\,000.

\begin{cuadrosolucion}
  \textbf{Solución:}
  \begin{enumerate}
    \item Función: $F(x) = 300x + 2\,000$ (afín, con costo fijo $b = 2\,000$ y precio por hora $m = 300$).
    \item Buscamos $x$ tal que $F(x) = 5\,000$:
          $$300x + 2\,000 = 5\,000 \Rightarrow 300x = 3\,000 \Rightarrow x = 10.$$
    \item Con \$5\,000 se puede estacionar 10 horas. Comprobación: $F(10) = 300 \cdot 10 + 2\,000 = 5\,000$ $\checkmark$
  \end{enumerate}
\end{cuadrosolucion}

\subsection{Ejercicios propuestos}

\begin{enumerate}[leftmargin=*, resume]
  \item \facil{} 3 kilos de manzanas cuestan \$2\,400. Escribe la función del costo en función de los kilos y calcula cuánto cuestan 5 kilos.

  \item \facil{} Un taxi cobra \$1\,000 de bajada de bandera y \$600 por kilómetro. Escribe la función del costo y calcula el costo de un viaje de 6 km.

  \item \medio{} Una fotocopiadora cobra \$50 por copia. Escribe la función del costo y calcula cuántas copias se pueden sacar con \$1\,750.

  \item \medio{} Un gimnasio cobra \$8\,000 de inscripción y \$5\,000 mensuales. Escribe la función del costo total en función de los meses y calcula el costo de 6 meses.

  \item \medio{} Con \$11\,000 en el estacionamiento del Ejemplo R3, ¿cuántas horas se pueden estacionar?

  \item \dificil{} Un vendedor gana \$250\,000 fijos más el 10\% de sus ventas. Escribe la función de su sueldo en función de las ventas y calcula cuánto debe vender para ganar \$400\,000.

\end{enumerate}
\newpage

% ============================================================
\section{Clave de Respuestas}
% ============================================================

\subsection*{Sección 1: Conceptos Básicos (Ejercicios 1--7)}

\begin{enumerate}[leftmargin=*, label=\textbf{\arabic*.}]
  \item a) afín \quad b) lineal \quad c) constante \quad d) lineal

  \item a) $m = 6$, $b = -1$ \quad b) $m = -1$, $b = 4$ \quad c) $m = \dfrac{3}{4}$, $b = 0$

  \item $F(0) = 5$, $F(1) = 7$, $F(-2) = 1$

  \item $G(2) = -5$, $G(-1) = 4$

  \item La pendiente es $m = 0$; la gráfica de $F(x) = mx$ pasa por el origen $(0, 0)$.

  \item $V(t) = 12\,000 - 800t$: es afín (hay un valor inicial fijo).

  \item $F(1) = m + b = 5$ y $F(3) = 3m + b = 11$. Restando: $2m = 6 \Rightarrow m = 3$, y $b = 2$. $F(x) = 3x + 2$.
\end{enumerate}

\newpage

\subsection*{Sección 2: Plano Cartesiano (Ejercicios 8--13)}

\begin{enumerate}[leftmargin=*, label=\textbf{\arabic*.}]
  \setcounter{enumi}{7}
  \item a) I \quad b) II \quad c) III \quad d) eje $Y$ \quad e) eje $X$

  \item a) $(3, 2)$ \quad b) $(-4, -1)$

  \item $x<0$, $y<0$: cuadrante III. $x>0$, $y<0$: cuadrante IV.

  \item Ubicación en el plano (ver figura del ejercicio resuelto R1).

  \item a) III \quad b) IV \quad c) III \quad d) eje $Y$ \quad e) II \quad f) eje $X$ \quad g) eje $Y$ \quad h) III

  \item Ejemplos: $(-3, 2)$ y $(-5, 4)$ (ambos con $x<0$, $y>0$ y suma $-1$).
\end{enumerate}

\newpage

\subsection*{Sección 3: Identificación del Tipo (Ejercicios 14--18)}

\begin{enumerate}[leftmargin=*, label=\textbf{\arabic*.}]
  \setcounter{enumi}{13}
  \item a) afín \quad b) lineal \quad c) constante ($F(x) = 6$) \quad d) lineal

  \item a) constante \quad b) creciente \quad c) decreciente

  \item Su pendiente es $m = 0$: es una función constante.

  \item Ejemplo: $F(x) = -2x + 5$.

  \item $m = \dfrac{8 - 2}{3 - 1} = 3$; luego $F(1) = 3 + b = 2 \Rightarrow b = -1$. $F(x) = 3x - 1$.
\end{enumerate}

\newpage

\subsection*{Sección 4: Gráficas (Ejercicios 19--24)}

\begin{enumerate}[leftmargin=*, label=\textbf{\arabic*.}]
  \setcounter{enumi}{18}
  \item Tabla: $x: -2,-1,0,1,2$; $F(x): 1,2,3,4,5$. Recta creciente que pasa por $(0,3)$.

  \item Tabla: $x: -2,-1,0,1,2$; $F(x): 5,3,1,-1,-3$. Recta decreciente que pasa por $(0,1)$.

  \item a) Cortes: $(0,-6)$ y $(2,0)$. \quad b) Cortes: $(0,8)$ y $(4,0)$.

  \item a) Cortes: $(0,4)$ y $(3,0)$. \quad b) Cortes: $(0,-6)$ y $(9,0)$.

  \item a) $F(x) = -3x$: lineal (pasa por el origen), decreciente. \quad b) $F(x) = -6$: constante (recta horizontal $y = -6$).

  \item $F(x) = 0 \Rightarrow \dfrac{2}{3}x = 6 \Rightarrow x = 9 > 0$: corta el eje $X$ en un valor positivo.
\end{enumerate}

\newpage

\subsection*{Sección 5: Pendiente (Ejercicios 25--30)}

\begin{enumerate}[leftmargin=*, label=\textbf{\arabic*.}]
  \setcounter{enumi}{24}
  \item a) creciente \quad b) decreciente \quad c) decreciente \quad d) constante

  \item a) $m = 3$ \quad b) $m = 3$ \quad c) $m = -2$

  \item a) $m = \dfrac{-3 - 3}{2 - (-1)} = \dfrac{-6}{3} = -2$ \quad b) $m = \dfrac{9 - 1}{4 - 4} = \dfrac{8}{0}$: no está definida (recta vertical; ¡no es función!).

  \item $m = \dfrac{2}{3} > 0$: asciende.

  \item $F(x) = -2x + 6$ (verifica: $F(1) = -2 + 6 = 4$). $\checkmark$

  \item $m = \dfrac{9 - 5}{(a + 2) - a} = \dfrac{4}{2} = 2$.
\end{enumerate}

\newpage

\subsection*{Sección 6: Problemas de Aplicación (Ejercicios 31--37)}

\begin{enumerate}[leftmargin=*, label=\textbf{\arabic*.}]
  \setcounter{enumi}{30}
  \item $m = \dfrac{2400}{3} = 800$; $F(x) = 800x$; $F(5) = 4\,000$.

  \item $F(x) = 600x + 1\,000$; $F(6) = 4\,600$.

  \item $F(x) = 50x$; $50x = 1\,750 \Rightarrow x = 35$ copias.

  \item $F(x) = 5\,000x + 8\,000$; $F(6) = 38\,000$.

  \item $300x + 2\,000 = 11\,000 \Rightarrow x = 30$ horas.

  \item Sueldo $= 0{,}1v + 250\,000$; $0{,}1v + 250\,000 = 400\,000 \Rightarrow v = 1\,500\,000$ en ventas.
\end{enumerate}

\end{document}
